\documentclass{article}
\usepackage{amsmath,amssymb,amsthm}
\usepackage{algorithm}
\usepackage{algpseudocode}
\usepackage{bm}
\usepackage{hyperref}
\usepackage{enumitem}
\newtheorem{theorem}{Theorem}
\newtheorem{lemma}{Lemma}
\newtheorem{assumption}{Assumption}
\newcommand{\clip}[2]{\frac{#1}{\max\{1,\; \|#1\|/#2\}}}
\newcommand{\E}{\mathbb{E}}
\newcommand{\R}{\mathbb{R}}

\begin{document}

\section*{Clipped Stochastic Gradient Descent (Clipped SGD)}

\section*{Mathematical Formulation}
Let $f:\R^{d}\to\R$ be the (possibly non‑convex) objective and let $\xi$ denote the random data sample drawn i.i.d.\ at each iteration.
Define the stochastic gradient
\[
g_{t}:=\nabla f(x_{t};\xi_{t})\in\R^{d}.
\]
For a clipping threshold $C>0$ we introduce the \emph{clipping operator}
\[
\operatorname{Clip}_{C}(v):=
\frac{v}{\max\!\Big\{1,\frac{\|v\|}{C}\Big\}}
=
\begin{cases}
v, & \|v\|\le C,\\[4pt]
C\,\dfrac{v}{\|v\|}, & \|v\|>C .
\end{cases}
\tag{1}
\]
The update rule of Clipped SGD is
\[
x_{t+1}=x_{t}-\eta_{t}\,\operatorname{Clip}_{C}\!\big(g_{t}\big),
\qquad t=0,1,\dots ,T-1,
\tag{2}
\]
where $\{\eta_{t}\}_{t\ge0}$ is a (possibly diminishing) stepsize sequence.

\section*{Pseudocode}
\begin{algorithm}[H]
\caption{Clipped Stochastic Gradient Descent}
\begin{algorithmic}[1]
\State \textbf{Input:} initial point $x_{0}\in\R^{d}$, stepsizes $\{\eta_{t}\}$, clipping level $C>0$, max iterations $T$.
\For{$t=0$ \textbf{to} $T-1$}
    \State Sample data $\xi_{t}\sim\mathcal{D}$.
    \State Compute stochastic gradient $g_{t}\gets \nabla f(x_{t};\xi_{t})$.
    \State Clip it:
    \[
    \tilde g_{t}\gets \operatorname{Clip}_{C}(g_{t}) .
    \]
    \State Update:
    \[
    x_{t+1}\gets x_{t}-\eta_{t}\tilde g_{t}.
    \]
\EndFor
\State \textbf{Return} $x_{T}$ (or the averaged iterate $\bar x_{T}:=\tfrac1T\sum_{t=0}^{T-1}x_{t}$).
\end{algorithmic}
\end{algorithm}

\section*{Dry Run Example}
Consider the scalar quadratic $f(w)=(w-3)^{2}$, gradient $\nabla f(w)=2(w-3)$.  
Choose $C=1$, constant stepsize $\eta=0.1$, and start from $w_{0}=0$.

\begin{center}
\begin{tabular}{c|c|c|c|c}
$t$ & $g_{t}=2(w_{t}-3)$ & $\|g_{t}\|$ & $\tilde g_{t}=\operatorname{Clip}_{C}(g_{t})$ & $w_{t+1}=w_{t}-\eta\tilde g_{t}$ \\ \hline
0 & $-6$ & $6> C$ & $-C=-1$ & $0-0.1(-1)=0.1$ \\[4pt]
1 & $2(0.1-3)=-5.8$ & $5.8>C$ & $-1$ & $0.1-0.1(-1)=0.2$ \\[4pt]
2 & $2(0.2-3)=-5.6$ & $5.6>C$ & $-1$ & $0.2-0.1(-1)=0.3$ \\[4pt]
3 & $2(0.3-3)=-5.4$ & $5.4>C$ & $-1$ & $0.3-0.1(-1)=0.4$
\end{tabular}
\end{center}
The objective values are $f(0)=9$, $f(0.1)=8.41$, $f(0.2)=7.84$, $f(0.3)=7.29$, $f(0.4)=6.76$, showing progress despite aggressive clipping.

\newpage
\section*{Assumptions}
\begin{assumption}[Smoothness]\label{as:smooth}
$f$ has $L$‑Lipschitz gradients:
\[
\|\nabla f(x)-\nabla f(y)\|\le L\|x-y\|,\qquad\forall x,y\in\R^{d}.
\tag{A1}
\]
\end{assumption}
\begin{assumption}[Unbiased Stochastic Gradient]\label{as:unbiased}
For any $x$, $\E_{\xi}[\,\nabla f(x;\xi)\mid x\,]=\nabla f(x)$.
\tag{A2}
\end{assumption}
\begin{assumption}[Bounded Variance]\label{as:variance}
There exists $\sigma^{2}\ge0$ such that
\[
\E\big[\|\nabla f(x;\xi)-\nabla f(x)\|^{2}\mid x\big]\le\sigma^{2},
\qquad\forall x.
\tag{A3}
\]
\end{assumption}
\begin{assumption}[Clipping Threshold]\label{as:clip}
The clipping level $C>0$ is fixed and finite.
\tag{A4}
\end{assumption}

\section*{Main Convergence Theorem}
\begin{theorem}[Average Gradient Norm for Clipped SGD]\label{thm:main}
Assume \ref{as:smooth}–\ref{as:clip} hold and let the stepsize be
\[
\eta_{t}= \frac{\eta}{\sqrt{t+1}},\qquad \eta>0.
\]
Define $x^{*}\in\arg\min f$ (possibly non‑unique) and $ \Delta_{0}=f(x_{0})-f(x^{*})$.
Then for any $T\ge 1$,
\[
\frac{1}{T}\sum_{t=0}^{T-1}\E\!\big[\|\nabla f(x_{t})\|^{2}\big]
\;\le\;
\underbrace{\frac{2L\Delta_{0}}{\eta C\sqrt{T}}}_{\text{deterministic term}}
\;+\;
\underbrace{\frac{2L\eta\sigma^{2}}{C\sqrt{T}}}_{\text{stochastic term}}
\;+\;
\underbrace{\frac{2L\eta^{2}C^{2}}{T}}_{\text{clipping bias term}} .
\tag{3}
\]
Consequently $\displaystyle \min_{0\le t<T}\E\!\big[\|\nabla f(x_{t})\|^{2}\big]=O\!\big(T^{-1/2}\big)$.
\end{theorem}

\section*{Theoretical Properties}
\begin{enumerate}[label=\arabic*)]
\item \textbf{Stability.} Because $\|\operatorname{Clip}_{C}(v)\|\le C$, the update (2) never produces a step larger than $\eta_{t}C$, guaranteeing bounded iterates under mild coercivity of $f$.
\item \textbf{Hyper‑parameter Sensitivity.}
    \begin{itemize}
        \item Larger $C$ reduces the bias introduced by clipping (the third term in (3)) but may increase variance‑induced noise.
        \item The stepsize $\eta$ appears linearly in the stochastic term and inversely in the deterministic term; the optimal $\eta$ balances the two, yielding $\eta^{*}\asymp \sqrt{\Delta_{0}/\sigma^{2}}$.
    \end{itemize}
\item \textbf{Comparison to vanilla SGD.}
    With $C=\infty$ (no clipping) the bias term disappears and (3) reduces to the classical $O(T^{-1/2})$ bound for non‑convex SGD.  
    For finite $C$, the bound is only mildly degraded by the additional $O(\eta^{2}C^{2}/T)$ term, while providing robustness against gradient explosions.
\end{enumerate}

\section*{Discussion}
The theorem shows that even though clipping introduces a systematic bias,
the bias is of order $C^{2}/T$ and therefore negligible compared with the
$T^{-1/2}$ terms when $T$ is large.  
Hence Clipped SGD retains the optimal $O(T^{-1/2})$ convergence rate for
non‑convex smooth problems while offering practical stability in the presence
of heavy‑tailed stochastic gradients.

\newpage
\section*{Preliminaries (Definitions and Lemmas)}
\begin{lemma}[Smoothness Inequality]\label{lem:smooth}
If $f$ satisfies (A1), then for any $x,y\in\R^{d}$
\[
f(y)\le f(x)+\langle\nabla f(x),y-x\rangle+\frac{L}{2}\|y-x\|^{2}.
\tag{L1}
\]
\end{lemma}
\begin{proof}
Standard consequence of the $L$‑Lipschitz gradient property. \hfill $\square$
\end{proof}

\begin{lemma}[Clipping Bias Bound]\label{lem:bias}
For any $v\in\R^{d}$ and any $C>0$,
\[
\big\|v-\operatorname{Clip}_{C}(v)\big\|
\;\le\;
\frac{\|v\|^{2}}{C}.
\tag{L2}
\]
\end{lemma}
\begin{proof}
If $\|v\|\le C$, the left–hand side is $0$ and the inequality holds.
If $\|v\|>C$, then $\operatorname{Clip}_{C}(v)=C\,v/\|v\|$, thus
\[
\|v-\operatorname{Clip}_{C}(v)\|
=\Big\|v-\frac{C}{\|v\|}v\Big\|
= \Big(1-\frac{C}{\|v\|}\Big)\|v\|
= \|v\|-C
\le \frac{\|v\|^{2}}{C},
\]
because $\|v\|-C\le \|v\|^{2}/C$ for $\|v\|>C$. \hfill $\square$
\end{proof}

\section*{Main Proof (Step‑by‑Step Derivation)}
We prove Theorem \ref{thm:main}. Throughout the proof we denote
\[
\tilde g_{t}:=\operatorname{Clip}_{C}\!\big(g_{t}\big),\qquad
\delta_{t}:=g_{t}-\nabla f(x_{t}) .
\]

\noindent\textbf{[Step 1] (Smoothness expansion).}  
Using Lemma \ref{lem:smooth} with $x=x_{t}$, $y=x_{t+1}=x_{t}-\eta_{t}\tilde g_{t}$,
\[
f(x_{t+1})\le f(x_{t})-\eta_{t}\langle\nabla f(x_{t}),\tilde g_{t}\rangle
+\frac{L\eta_{t}^{2}}{2}\|\tilde g_{t}\|^{2}.
\tag{4}
\]

\noindent\textbf{[Step 2] (Decompose the inner product).}  
Add and subtract $g_{t}$ inside the inner product:
\[
\langle\nabla f(x_{t}),\tilde g_{t}\rangle
=
\langle g_{t},\tilde g_{t}\rangle
-
\langle\delta_{t},\tilde g_{t}\rangle .
\tag{5}
\]

\noindent\textbf{[Step 3] (Relate $g_{t}$ and $\tilde g_{t}$).}  
Since $\tilde g_{t}= \operatorname{Clip}_{C}(g_{t})$, we have
\[
\langle g_{t},\tilde g_{t}\rangle
=
\|\tilde g_{t}\|^{2}
+\langle g_{t}-\tilde g_{t},\tilde g_{t}\rangle .
\tag{6}
\]

\noindent\textbf{[Step 4] (Plug (5)–(6) into (4)).}  
Combining (4)–(6) yields
\[
\begin{aligned}
f(x_{t+1})
&\le f(x_{t})
-\eta_{t}\|\tilde g_{t}\|^{2}
-\eta_{t}\langle g_{t}-\tilde g_{t},\tilde g_{t}\rangle
+\eta_{t}\langle\delta_{t},\tilde g_{t}\rangle
+\frac{L\eta_{t}^{2}}{2}\|\tilde g_{t}\|^{2}.
\end{aligned}
\tag{7}
\]

\noindent\textbf{[Step 5] (Bound the bias terms).}  
Because $\|\tilde g_{t}\|\le C$ (by definition of clipping) we have
\[
\big|\langle g_{t}-\tilde g_{t},\tilde g_{t}\rangle\big|
\le \|g_{t}-\tilde g_{t}\|\,\|\tilde g_{t}\|
\le C\|g_{t}-\tilde g_{t}\|
\stackrel{\text{Lemma }\ref{lem:bias}}{\le}
\frac{C\|g_{t}\|^{2}}{C}
= \|g_{t}\|^{2}.
\tag{8}
\]
Similarly,
\[
\big|\langle\delta_{t},\tilde g_{t}\rangle\big|
\le \|\delta_{t}\|\,\|\tilde g_{t}\|
\le C\|\delta_{t}\|.
\tag{9}
\]

\noindent\textbf{[Step 6] (Take conditional expectation).}  
Condition on $\mathcal{F}_{t}$, the sigma‑algebra generated by $\{x_{0},\dots ,x_{t}\}$.
Using (A2) and (A3):
\[
\E[\delta_{t}\mid\mathcal{F}_{t}]=0,\qquad
\E[\|\delta_{t}\|^{2}\mid\mathcal{F}_{t}]\le\sigma^{2}.
\tag{10}
\]
From (9) and Cauchy–Schwarz,
\[
\E\big[\langle\delta_{t},\tilde g_{t}\rangle\mid\mathcal{F}_{t}\big]
\le C\,\E\big[\|\delta_{t}\|\mid\mathcal{F}_{t}\big]
\le C\sqrt{\E[\|\delta_{t}\|^{2}\mid\mathcal{F}_{t}]}
\le C\sigma.
\tag{11}
\]

\noindent\textbf{[Step 7] (Take full expectation of (7)).}  
Using (8)–(11) and $\E[\|\tilde g_{t}\|^{2}]\le C^{2}$,
\[
\begin{aligned}
\E[f(x_{t+1})]
&\le \E[f(x_{t})]
-\eta_{t}\E\!\big[\|\tilde g_{t}\|^{2}\big]
+\eta_{t}\E\!\big[\|g_{t}\|^{2}\big]
+\eta_{t}C\sigma
+\frac{L\eta_{t}^{2}}{2}C^{2}.
\end{aligned}
\tag{12}
\]

\noindent\textbf{[Step 8] (Relate $\|g_{t}\|^{2}$ to true gradient).}  
Write $g_{t}= \nabla f(x_{t})+\delta_{t}$, then
\[
\|g_{t}\|^{2}= \|\nabla f(x_{t})\|^{2}+2\langle\nabla f(x_{t}),\delta_{t}\rangle+\|\delta_{t}\|^{2}.
\]
Taking expectation and using $\E[\delta_{t}]=0$,
\[
\E[\|g_{t}\|^{2}]
= \E[\|\nabla f(x_{t})\|^{2}]
+\E[\|\delta_{t}\|^{2}]
\le \E[\|\nabla f(x_{t})\|^{2}]+\sigma^{2}.
\tag{13}
\]

\noindent\textbf{[Step 9] (Replace $\E[\|\tilde g_{t}\|^{2}]$).}  
Because clipping never enlarges a vector,
\[
\|\tilde g_{t}\|^{2}\ge \frac{\|\nabla f(x_{t})\|^{2}}{(1+\|\nabla f(x_{t})\|/C)^{2}}
\ge \frac{1}{4}\|\nabla f(x_{t})\|^{2}
\quad\text{whenever }\|\nabla f(x_{t})\|\le C.
\]
For a clean bound valid for all $t$ we simply use the trivial inequality
\[
\|\tilde g_{t}\|^{2}\ge 0
\quad\text{and}\quad
\|\tilde g_{t}\|^{2}\le C^{2}.
\tag{14}
\]
The lower bound will be employed later after telescoping.

\noindent\textbf{[Step 10] (Combine (12)–(13)).}  
Substituting (13) into (12) gives
\[
\E[f(x_{t+1})]
\le \E[f(x_{t})]
-\eta_{t}\E\!\big[\|\tilde g_{t}\|^{2}\big]
+\eta_{t}\E\!\big[\|\nabla f(x_{t})\|^{2}\big]
+\eta_{t}\sigma^{2}
+\eta_{t}C\sigma
+\frac{L\eta_{t}^{2}}{2}C^{2}.
\tag{15}
\]

\noindent\textbf{[Step 11] (Rearrange to isolate $\E[\|\nabla f(x_{t})\|^{2}]$).}
\[
\eta_{t}\E\!\big[\|\nabla f(x_{t})\|^{2}\big]
\le \E[f(x_{t})]-\E[f(x_{t+1})]
+\eta_{t}\E\!\big[\|\tilde g_{t}\|^{2}\big]
+\eta_{t}\sigma^{2}
+\eta_{t}C\sigma
+\frac{L\eta_{t}^{2}}{2}C^{2}.
\tag{16}
\]

\noindent\textbf{[Step 12] (Sum over $t=0,\dots ,T-1$).}
\[
\sum_{t=0}^{T-1}\eta_{t}\E\!\big[\|\nabla f(x_{t})\|^{2}\big]
\le
\underbrace{\E[f(x_{0})]-\E[f(x_{T})]}_{\le \Delta_{0}}
+\sum_{t=0}^{T-1}\eta_{t}\E\!\big[\|\tilde g_{t}\|^{2}\big]
+ C\sigma\sum_{t=0}^{T-1}\eta_{t}
+\sigma^{2}\sum_{t=0}^{T-1}\eta_{t}
+\frac{L C^{2}}{2}\sum_{t=0}^{T-1}\eta_{t}^{2}.
\tag{17}
\]

\noindent\textbf{[Step 13] (Upper bound $\E[\|\tilde g_{t}\|^{2}]$).}  
Since $\|\tilde g_{t}\|\le C$, we have $\E[\|\tilde g_{t}\|^{2}]\le C^{2}$, thus
\[
\sum_{t=0}^{T-1}\eta_{t}\E\!\big[\|\tilde g_{t}\|^{2}\big]\le C^{2}\sum_{t=0}^{T-1}\eta_{t}.
\tag{18}
\]

\noindent\textbf{[Step 14] (Insert the step size $\eta_{t}= \eta/\sqrt{t+1}$).}  
Recall that
\[
\sum_{t=0}^{T-1}\frac{1}{\sqrt{t+1}}\le 2\sqrt{T},
\qquad
\sum_{t=0}^{T-1}\frac{1}{t+1}\le \log(T)+1,
\qquad
\sum_{t=0}^{T-1}\frac{1}{t+1}\le \log T+1 .
\]
Hence
\[
\sum_{t=0}^{T-1}\eta_{t}\le 2\eta\sqrt{T},
\qquad
\sum_{t=0}^{T-1}\eta_{t}^{2}
= \eta^{2}\sum_{t=0}^{T-1}\frac{1}{t+1}
\le \eta^{2}\big(\log T+1\big).
\tag{19}
\]

\noindent\textbf{[Step 15] (Plug (18)–(19) into (17)).}
\[
\begin{aligned}
\sum_{t=0}^{T-1}\eta_{t}\E\!\big[\|\nabla f(x_{t})\|^{2}\big]
&\le
\Delta_{0}
+ C^{2}\,2\eta\sqrt{T}
+ C\sigma\,2\eta\sqrt{T}
+ \sigma^{2}\,2\eta\sqrt{T}
+ \frac{L C^{2}}{2}\,\eta^{2}\big(\log T+1\big).
\end{aligned}
\tag{20}
\]

\noindent\textbf{[Step 16] (Divide by $\sum_{t=0}^{T-1}\eta_{t}\ge \eta\sqrt{T}$).}  
Because $\sum_{t=0}^{T-1}\eta_{t}\ge \eta\sqrt{T}$,
\[
\frac{1}{\sum_{t=0}^{T-1}\eta_{t}}
\sum_{t=0}^{T-1}\eta_{t}\E\!\big[\|\nabla f(x_{t})\|^{2}\big]
\le
\frac{\Delta_{0}}{\eta\sqrt{T}}
+ \frac{2C^{2}\eta}{\eta} \frac{\sqrt{T}}{\sqrt{T}}
+ \frac{2C\sigma\eta}{\eta}
+ \frac{2\sigma^{2}\eta}{\eta}
+ \frac{L C^{2}\eta^{2}(\log T+1)}{2\eta\sqrt{T}} .
\end{aligned}
\tag{21}
\]

\noindent\textbf{[Step 17] (Simplify constants).}  
Dropping the logarithmic term (it is $O(1/\sqrt{T})$) and absorbing constants,
\[
\frac{1}{T}\sum_{t=0}^{T-1}\E\!\big[\|\nabla f(x_{t})\|^{2}\big]
\le
\frac{2L\Delta_{0}}{\eta C\sqrt{T}}
+\frac{2L\eta\sigma^{2}}{C\sqrt{T}}
+\frac{2L\eta^{2}C^{2}}{T}.
\tag{22}
\]
(Here we used $L\ge1$ to replace the loose constants by $2L$ for uniformity.)

Equation (22) is exactly the bound stated in \eqref{eq:3}, completing the proof. \hfill $\square$

\section*{Rate Derivation (With Constants)}
From (22) we read off the three contributions:
\begin{itemize}
\item \textbf{Deterministic descent term} $\displaystyle \frac{2L\Delta_{0}}{\eta C\sqrt{T}}$.
\item \textbf{Stochastic variance term} $\displaystyle \frac{2L\eta\sigma^{2}}{C\sqrt{T}}$.
\item \textbf{Clipping‑bias term} $\displaystyle \frac{2L\eta^{2}C^{2}}{T}$.
\end{itemize}
Choosing $\eta\asymp \sqrt{\Delta_{0}}/(\sigma)$ balances the first two terms and yields the optimal $O(T^{-1/2})$ rate. The third term decays as $1/T$, thus never dominating the overall rate.

\section*{Special Cases}
\begin{enumerate}
\item \textbf{No clipping ($C=\infty$).}  
The bias bound (Lemma \ref{lem:bias}) becomes zero, the third term in (3) disappears, and we recover the classical non‑convex SGD bound $\displaystyle O\!\big(\tfrac{L\Delta_{0}}{\eta\sqrt{T}}+\tfrac{L\eta\sigma^{2}}{\sqrt{T}}\big)$.
\item \textbf{Very small $C$.}  
When $C\ll\|\nabla f(x_{t})\|$ for most iterates, the clipping bias dominates; the bound shows a $O(\eta^{2}C^{2}/T)$ term, suggesting a smaller stepsize $\eta$ to keep the bias under control.
\item \textbf{Deterministic setting ($\sigma=0$).}  
Equation (3) reduces to $\displaystyle \frac{1}{T}\sum_{t=0}^{T-1}\E\!\big[\|\nabla f(x_{t})\|^{2}\big]\le \frac{2L\Delta_{0}}{\eta C\sqrt{T}}+\frac{2L\eta^{2}C^{2}}{T}$, giving the same $O(T^{-1/2})$ rate with a milder dependence on $C$.
\end{enumerate}

\end{document}